\subsection{Structural Summary}

The preceding sections separate two ideas that are often confused:

\begin{itemize}
    \item Python names are dynamically bound.
    \item Python objects are strongly typed at runtime.
\end{itemize}

Python is flexible at assignment and strict at operation. A name may point to different objects at different moments, but each object brings its own type rules through \texttt{ob\_type}.

\subsubsection{Categorizing Python as a Dynamically Bound, Strongly Verified Runtime Typing Environment}

Python is dynamically typed because names have no fixed declared types. There is no C-style \texttt{int value;} declaration. Assignment binds a name to an object; later rebinding to another type is legal.

Python is strongly verified because the interpreter consults object types when operations execute. \texttt{"3" + 4} fails not because the name is wrong, but because the string object's addition protocol rejects an integer operand.

The precise classification:

\begin{quote}
Python uses dynamically bound names and strongly verified runtime operations.
\end{quote}

That implies:

\begin{itemize}
    \item names can be rebound freely across singular domains;
    \item objects carry runtime types via \texttt{ob\_type};
    \item operators consult object behavior through special methods;
    \item unsupported combinations fail with \texttt{TypeError} rather than silent conversion;
    \item explicit conversion is the programmer's responsibility at domain boundaries.
\end{itemize}

\subsubsection{The Practical Rule: Names Are Flexible, Objects Keep Their Runtime Type}

A name may change what it references, but the object it currently references keeps its type and behavior. The name \texttt{thing} does not carry \texttt{bit\_length} or \texttt{upper}; the current object does.

For \texttt{+}, integer addition and string concatenation are different dispatches on the same token. The object type decides which \texttt{\_\_add\_\_} implementation runs. If the operation is undefined, Python rejects it.

The structural summary for this chapter:

\begin{verbatim}
assignment:    flexible name-to-object binding
reassignment:  same name, possibly different object type
type():        exact current runtime type
isinstance():  category or family test (mind bool/int)
operation:     checked against current objects
conversion:    explicit when crossing type boundaries
\end{verbatim}

The final rule:

\begin{quote}
Python is permissive about rebinding names, but strict about what objects can do together.
\end{quote}
